% Copyright (c) 2026 Frankie Feng-Cheng WANG. All rights reserved.
% Repository: https://github.com/FrankieeW/formalising-mathematics-notes

%----------------------------------------------------------------------------------------
% Coursework/Project Requirements
% What is the mark scheme?
% Here is the thought experiment.
% Consider an external examiner (a professor from another UK
% university).
% They have heard of Lean and that know it’s a “theorem prover”.
% They have no idea how it works.
% What would they think of your project?
% Would they be able to understand something about what you
% are doing?
% Would they think “this is pretty cool, we should be teaching this
% at my university”?
%
% For the first project. I am looking for:
% ▶ some Lean code (150–200 lines of code including
% comments? More if you like?)
% ▶ A pdf explaining what is going on (5 pages? More if you
% like?)
% You won’t get penalised for doing too much (but it will take
% longer).
% I’m expecting about 15 (or more) hours of work.
%----------------------------------------------------------------------------------------
% PACKAGES AND OTHER DOCUMENT CONFIGURATIONS
%----------------------------------------------------------------------------------------

\documentclass[
  % UTF8, % Encoding %not needed
  12pt, % Default font size, values between 10pt-12pt are allowed
  % letterpaper, % Uncomment for US letter paper size, the default paper size is a4paper
  % Chinese, % Uncomment for Chinese document
]{assignment}

\usepackage{amsthm} % For proof environment
\newtheorem{lemma}{Lemma} % define lemma environment

% Definition environment matching problem environment style from assignment.cls
% Using framed environment like the problem environment in assignment.cls
\newenvironment{mydef} {
  \begin{framed}
} {
  \end{framed}
  \bigskip
}

% \usepackage{pdfpages} % 用于插入 PDF
\usepackage{fontspec}
\setmonofont{FreeMono}
\usepackage{hyperref} % ← 必须显式、提前
\usepackage{minted}
\newmintinline[lean]{lean4}{bgcolor=white}
\newminted[leancode]{lean4}{fontsize=\footnotesize,breaklines,breakanywhere,tabsize=4,showspaces=false}
\newcommand{\leancodeurl}{}
\newcommand{\setleancodeurl}[1]{\renewcommand{\leancodeurl}{#1}}
\renewcommand{\theFancyVerbLine}{\arabic{FancyVerbLine}}
\renewcommand{\FancyVerbFormatLine}[1]{\href{\leancodeurl\#L\arabic{FancyVerbLine}}{#1}}

% Original environment for inline code (deprecated, use leancodefile instead)
\newenvironment{leancodegit}[3][]{\VerbatimEnvironment\setleancodeurl{#3}\begin{leancode}[#1]}{\end{leancode}}

% New command to import code directly from .lean files
% Usage: \leancodefile[minted options]{local/path/to/file.lean}{github_url}
% Example: \leancodefile[firstline=20,lastline=25]{../lean/GroupAction/Defs.lean}{https://github.com/.../Defs.lean}
\newcommand{\leancodefile}[3][]{%
  \setleancodeurl{#3}%
  \inputminted[fontsize=\footnotesize,breaklines,breakanywhere,tabsize=4,showspaces=false,linenos=true,#1]{lean4}{#2}%
}

\usemintedstyle{tango}
% \usepackage{hyperref}
% \usepackage{embedfile}
% \usepackage{appendix}
\usepackage{mathtools}
\usepackage{amssymb}
\usepackage{xcolor}
\newcommand\nb{\addtocounter{equation}{1}\tag{\theequation}}
\pgfplotsset{compat=1.18}
% \renewcommand{\thesection}{Exercise \arabic{section}}


%-------------------------------------------------------------------------------
% ASSIGNMENT INFORMATION
%-------------------------------------------------------------------------------

% Fix for fancyhdr warning: increase headheight and adjust topmargin
\setlength{\headheight}{29.34845pt}
\addtolength{\topmargin}{-17.34845pt}

\title{Formalizing Group Action in Lean} % Assignment title
\author{Frankie Feng-Cheng WANG} % Student name
% \studentid{219046894} % Student id Leicester
\email{maths@frankie.wang} % Student email
\github{https://github.com/FrankieeW/GroupAction} % GitHub URL
\date{\today} % Due date
\institute{Department of Mathematics\\Imperial College London} % Institute or school name
\course{MATH70040-Formalising Mathematics} % Course or class name
\lecturer{Dr Bhavik Mehta} % Lecturer or teacher in charge of the assignment

%----------------------------------------------------------------------------------------
\begin{document}
\maketitle % Output the assignment title, created automatically using the information in the custom commands above
\tableofcontents

%----------------------------------------------------------------------------------------
% ASSIGNMENT CONTENT
%----------------------------------------------------------------------------------------
\section{Introduction}

This report presents a Lean~4 formalisation of \textbf{group actions}, organised around three interrelated themes:
\begin{enumerate}
	\item \textbf{Core definitions:} the \texttt{GroupAction} typeclass, together with the notions of faithfulness and transitivity;
	\item \textbf{Concrete examples:} a selection of standard group actions illustrating the abstract definitions;
	\item \textbf{Key theorems:} the construction of the permutation representation and the subgroup structure of stabilisers.
\end{enumerate}

The mathematical development and theorem numbering follow
Fraleigh and Katz, \textit{A First Course in Abstract Algebra},
Addison--Wesley, 2003, Section~16 (Group Actions).

The complete Lean source code is available on GitHub:

\url{https://github.com/FrankieeW/GroupAction}.

All code snippets in this report are hyperlinked to the corresponding lines in the repository.

The principal results formalised in this project are the existence of the permutation representation associated to a group action and the fact that stabilisers form subgroups. These are stated below and proved in later sections.

\begin{quote}
	\textbf{Theorem 16.3.}
	Let $X$ be a $G$-set. For each $g \in G$, the map
	\[
		\sigma_g : X \to X, \qquad \sigma_g(x) = g \cdot x,
	\]
	is a permutation of $X$. The assignment $\phi : G \to \mathrm{Sym}(X)$ defined by
	$\phi(g) = \sigma_g$ is a group homomorphism, and for all $g \in G$ and $x \in X$,
	\[
		\phi(g)(x) = g \cdot x.
	\]
\end{quote}

\begin{quote}
	\textbf{Theorem 16.12.}
	Let $X$ be a $G$-set and let $x \in X$. The stabiliser of $x$ in $G$, defined by
	\[
		G_x = \{\, g \in G \mid g \cdot x = x \,\},
	\]
	is a subgroup of $G$.
\end{quote}


\iffalse
	\section{Background and Design Choices}
	\subsection{Why a custom action class}
	Mathlib already provides \texttt{MulAction}, but I define a small custom class
	\texttt{GroupAction} with just two axioms: compatibility with multiplication and
	the identity action. This keeps the development minimal and makes it easy to map
	Lean lines to textbook lines. The cost is that I cannot use existing \texttt{MulAction}
	lemmas, but the benefit is transparency.

	\subsection{What Lean is checking}
	Lean treats each statement as a typed object. When I define \texttt{sigmaPerm}, Lean
	checks that the inverse function is a two-sided inverse. When I define \texttt{phi},
	Lean checks that its image lies in the group \texttt{Equiv.Perm X} and that
	multiplication is permutation composition. This explicit checking prevents hidden
	gaps that can slip into informal proofs.
\fi

\section{Definitions}
This section introduces the group action class and the standing assumptions used
throughout the file.
\subsection{Group Action}
\begin{mydef}
	A \textbf{group\ action} of a group $G$ on a set $X$ is a map
	\[
		\cdot : G \times X \to X,
	\]
	satisfying the following axioms for all $g_1, g_2 \in G$ and $x \in X$:
	\begin{itemize}
		\item (Compatibility) $(g_1 g_2) \cdot x = g_1 \cdot (g_2 \cdot x)$;
		\item (Identity) $1 \cdot x = x$.
	\end{itemize}
\end{mydef}
I first declare a minimal class for group actions:
\leancodefile[firstline=20,lastline=23,firstnumber=20]{../lean/GroupAction/Defs.lean}{https://github.com/FrankieeW/GroupAction/blob/v1.0.0/lean/GroupAction/Defs.lean}


Throughout the file I work with a group $G$ acting on a type $X$:
\leancodefile[firstline=13,lastline=13,firstnumber=13]{../lean/GroupAction/Basic.lean}{https://github.com/FrankieeW/GroupAction/blob/v1.0.0/lean/GroupAction/Basic.lean}

\section{Concrete Examples of Group Actions}
This section presents several concrete instances of the \texttt{GroupAction} class,
illustrating how the abstract definition applies to familiar mathematical structures.

\subsection{Symmetric group acting on a set}
Let $X$ be any type. The symmetric group $\mathrm{Sym}(X)$ acts on $X$ by evaluation:
for $\sigma \in \mathrm{Sym}(X)$ and $x \in X$, define $\sigma \cdot x := \sigma(x)$.

\leancodefile[firstline=19,lastline=26,firstnumber=19]{../lean/GroupAction/Examples.lean}{https://github.com/FrankieeW/GroupAction/blob/v1.0.0/lean/GroupAction/Examples.lean}

The axiom proofs are immediate by reflexivity: composition of permutations and function
evaluation commute definitionally in Lean.

Two standard properties of this action are faithful and transitive:
\leancodefile[firstline=28,lastline=42,firstnumber=28]{../lean/GroupAction/Examples.lean}{https://github.com/FrankieeW/GroupAction/blob/v1.0.0/lean/GroupAction/Examples.lean}
The instance \texttt{DecidableEq X} is required because \texttt{Equiv.swap} constructs a
permutation by branching on whether two elements are equal, which needs decidable equality.

\subsection{Group acting on itself by left multiplication}
Every group $G$ acts on itself by left multiplication: $g_1 \cdot g_2 := g_1 * g_2$.

\leancodefile[firstline=44,lastline=51,firstnumber=44]{../lean/GroupAction/Examples.lean}{https://github.com/FrankieeW/GroupAction/blob/v1.0.0/lean/GroupAction/Examples.lean}

The \texttt{ga\_mul} axiom follows from associativity of group multiplication, and
\texttt{ga\_one} from the identity axiom.

\subsection{Subgroup acting on the group}
A subgroup $H \leq G$ acts on $G$ by left multiplication. Elements of $H$ are coerced
to elements of $G$ using \texttt{(h : G)}.

\leancodefile[firstline=55,lastline=64,firstnumber=55]{../lean/GroupAction/Examples.lean}{https://github.com/FrankieeW/GroupAction/blob/v1.0.0/lean/GroupAction/Examples.lean}

The \texttt{simp} tactic handles the coercion and applies associativity and identity axioms.

\subsection{Conjugation action of a subgroup on itself}
A subgroup $H$ acts on itself by conjugation: $h_1 \cdot h_2 := h_1 h_2 h_1^{-1}$.

\leancodefile[firstline=66,lastline=84,firstnumber=66]{../lean/GroupAction/Examples.lean}{https://github.com/FrankieeW/GroupAction/blob/v1.0.0/lean/GroupAction/Examples.lean}
% mul_assoc.{u_1} {G : Type u_1} [Semigroup G] (a b c : G) : a * b * c = a * (b * c)
The \texttt{mul\_assoc} lemma is from mathlib and states associativity of multiplication:

\begin{minted}[fontsize=\footnotesize,breaklines,breakanywhere]{lean4}
mul_assoc.{u_1} {G : Type u_1} [Semigroup G] (a b c : G) :
a * b * c = a * (b * c)
\end{minted}



The \texttt{ga\_mul} proof can easily be done by \texttt{simp [mul\_assoc]}, but I expanded it to show the steps explicitly by using \texttt{rw}.



For example the \texttt{rw [← mul\_assoc g₁]} will aotumatically find the passible situation looks like \texttt{g₁* (b * c)}, because we only give one argument \texttt{g₁} to \texttt{← mul\_assoc}.
% g₁ * (g₂ * g₃ * g₂⁻¹)-> g₁ * (g₂ * g₃) * g₂⁻¹
In detail, it will rewrite the left hand side \texttt{g₁ * (g₂ * g₃ * g₂⁻¹)} to \texttt{g₁ * (g₂ * g₃) * g₂⁻¹} according to the right to left direction of the lemma \texttt{mul\_assoc}.
\subsection{Scalar action on complex vector spaces}
The multiplicative group $\mathbb{C}^\times$ of nonzero complex numbers acts on
$\mathbb{C}^n$ by scalar multiplication.

\leancodefile[firstline=88,lastline=98,firstnumber=88]{../lean/GroupAction/Examples.lean}{https://github.com/FrankieeW/GroupAction/blob/v1.0.0/lean/GroupAction/Examples.lean}

Here \texttt{ℂˣ} denotes the units of $\mathbb{C}$ (invertible elements), and
\texttt{Fin n → ℂ} represents functions from $\{0, \ldots, n-1\}$ to $\mathbb{C}$
(i.e., $n$-dimensional complex vectors). The \texttt{ext} tactic proves function equality
pointwise.

\subsection{Scalar action via units}
The unit group $\alpha^\times$ of a monoid $\alpha$ acts on $\alpha^n$ by scalar multiplication.

\leancodefile[firstline=103,lastline=115,firstnumber=103]{../lean/GroupAction/Examples.lean}{https://github.com/FrankieeW/GroupAction/blob/v1.0.0/lean/GroupAction/Examples.lean}

This generalizes the complex example by abstracting over the coefficient monoid.

\subsection{Dihedral group action on a square}
Let $D_4$ be the symmetry group of a square, acting on $\mathbb{Z}/4\mathbb{Z}$.

\leancodefile[firstline=134,lastline=155,firstnumber=134]{../lean/GroupAction/Examples.lean}{https://github.com/FrankieeW/GroupAction/blob/v1.0.0/lean/GroupAction/Examples.lean}

\section{Permutation Representation}
The core construction is the map $\sigma_g : X \to X$ given by $\sigma_g(x) = g \cdot x$.
The inverse of $\sigma_g$ is $\sigma_{g^{-1}}$, so $\sigma_g$ is a permutation.
This yields the permutation representation $\phi : G \to \mathrm{Sym}(X)$.

\subsection{Proof sketch (informal)}
To show that $\sigma_g$ is a bijection, compute
$\sigma_{g^{-1}}(\sigma_g(x)) = (g^{-1} g) \cdot x = 1 \cdot x = x$, and similarly
$\sigma_g(\sigma_{g^{-1}}(x)) = x$. The representation is defined by
$\phi(g) = \sigma_g$, and multiplicativity follows from the action axiom:
\[
	\phi(g_1 g_2)(x) = (g_1 g_2) \cdot x = g_1 \cdot (g_2 \cdot x) = (\phi(g_1)\phi(g_2))(x).
\]

\subsection{Lean code}
First, define the underlying action map:
\leancodefile[firstline=26,lastline=29,firstnumber=26]{../lean/GroupAction/Permutation.lean}{https://github.com/FrankieeW/GroupAction/blob/v1.0.0/lean/GroupAction/Permutation.lean}

Then package it as a permutation using the inverse action:
\leancodefile[firstline=31,lastline=63,firstnumber=31]{../lean/GroupAction/Permutation.lean}{https://github.com/FrankieeW/GroupAction/blob/v1.0.0/lean/GroupAction/Permutation.lean}
Using \texttt{\#check} on \texttt{sigmaPerm} confirms it has type
\leancodefile[firstline=64,lastline=64,firstnumber=64]{../lean/GroupAction/Permutation.lean}{https://github.com/FrankieeW/GroupAction/blob/v1.0.0/lean/GroupAction/Permutation.lean}
This confirms the map $g \mapsto \sigma_g$ is a well-defined function from $G$ to $\mathrm{Sym}(X)$.
Define the representation and its basic properties:
\leancodefile[firstline=66,lastline=99,firstnumber=66]{../lean/GroupAction/Permutation.lean}{https://github.com/FrankieeW/GroupAction/blob/v1.0.0/lean/GroupAction/Permutation.lean}

Finally, package the existence statement:
\leancodefile[firstline=101,lastline=107,firstnumber=101]{../lean/GroupAction/Permutation.lean}{https://github.com/FrankieeW/GroupAction/blob/v1.0.0/lean/GroupAction/Permutation.lean}

\section{Stabilizers}
For a fixed $x \in X$, the stabilizer is the subset
\[
	G_x = \{ g \in G \mid g \cdot x = x \}.
\]
In Lean this is first defined as a set, then shown to be the carrier of a subgroup,
and finally packaged as a \texttt{Subgroup}.

\subsection{Proof sketch (informal)}
The identity element fixes every $x$, so $1 \in G_x$. If $g_1$ and $g_2$ fix $x$, then
$(g_1 g_2)\cdot x = g_1 \cdot (g_2 \cdot x) = g_1 \cdot x = x$, so $g_1 g_2 \in G_x$.
If $g$ fixes $x$, then
$g^{-1} \cdot x = g^{-1} \cdot (g \cdot x) = (g^{-1} g) \cdot x = x$,
so $g^{-1} \in G_x$.

\subsection{Lean code}
Start from the set definition:
\leancodefile[firstline=22,lastline=25,firstnumber=22]{../lean/GroupAction/Stabilizer.lean}{https://github.com/FrankieeW/GroupAction/blob/v1.0.0/lean/GroupAction/Stabilizer.lean}

Package it as a subgroup by checking closure:
% 加一段介绍 carrier one_mem mul_mem inv_mem
Define the subgroup structure by providing proofs of the subgroup axioms:
\begin{itemize}
	\item \texttt{carrier}: the underlying set is \texttt{stabilizerSet x}.
	\item \texttt{one\_mem'}: the identity fixes every $x$ by the \texttt{ga\_one} axiom.
	\item \texttt{mul\_mem'}: if $g_1$ and $g_2$ fix $x$, then so does $g_1 g_2$ by the
	      \texttt{ga\_mul} axiom.
	\item \texttt{inv\_mem'}: if $g$ fixes $x$, then so does $g^{-1}$ by combining
	      \texttt{ga\_mul} and \texttt{ga\_one}.
\end{itemize}
\leancodefile[firstline=27,lastline=52,firstnumber=27]{../lean/GroupAction/Stabilizer.lean}{https://github.com/FrankieeW/GroupAction/blob/v1.0.0/lean/GroupAction/Stabilizer.lean}

Then show the set is the carrier of a subgroup:
\leancodefile[firstline=54,lastline=57,firstnumber=54]{../lean/GroupAction/Stabilizer.lean}{https://github.com/FrankieeW/GroupAction/blob/v1.0.0/lean/GroupAction/Stabilizer.lean}

\iffalse
	\section{Main Theorem}
	The Lean theorem \texttt{group\_action\_to\_perm\_representation} packages the permutation
	representation together with its key properties. The proof uses the lemmas from the
	previous section: \texttt{phi\_apply}, \texttt{phi\_mul}, and \texttt{phi\_one}.
\fi

\section{What Lean Guarantees}
Formalizing these constructions in Lean provides guarantees that informal proofs cannot match.
When Lean accepts a definition or theorem, it has verified:
\begin{itemize}
	\item \textbf{Types are correct.} The function \texttt{phi : G → Equiv.Perm X} has the claimed type.
	      Lean checks that \texttt{sigmaPerm g} is actually a permutation (a bijection), not just a function.
	      If we mistakenly defined a non-bijective map, Lean would reject it.
	\item \textbf{No hidden assumptions.} Every lemma explicitly lists its hypotheses. If a proof uses
	      associativity, the code must invoke \texttt{mul\_assoc} or the \texttt{ga\_mul} axiom.
	      There are no ``clearly'' or ``it follows that'' steps that skip justification.
	\item \textbf{Axioms match definitions.} The \texttt{GroupAction} class requires \texttt{ga\_mul}
	      and \texttt{ga\_one}. If we omitted \texttt{ga\_one}, the proof of \texttt{phi\_one} would fail
	      because Lean would have no way to show $1 \cdot x = x$.
	\item \textbf{Subgroup structure is verified.} When we package the stabilizer as a \texttt{Subgroup},
	      Lean checks that we provided proofs of closure under multiplication, inverses, and identity.
	      The type system prevents us from claiming ``the stabilizer is a subgroup'' without proving it.
\end{itemize}

These checks prevent common errors: mistaken claims about injectivity, forgotten hypotheses,
and gaps in subgroup proofs. An external examiner can trust that the Lean code genuinely establishes
the mathematical claims, not just that it compiles.

\section{Reflection}
Writing the proofs in Lean forced me to make every step explicit. In particular:
\begin{itemize}
	\item I had to show explicitly that $\sigma_{g^{-1}}$ is an inverse.
	\item I had to unfold the definition of permutation multiplication.
\end{itemize}
The resulting code is not shorter than a textbook proof, but it is more precise.
An external examiner can read the surrounding explanations and then see that each
Lean line matches a specific mathematical claim.


\subsection{Specific challenges encountered}
Several technical hurdles emerged during the formalization:
\begin{itemize}
	\item \textbf{Typeclass resolution:} Ensuring Lean could automatically find the \texttt{GroupAction}
	      instance in complex contexts required careful use of explicit type annotations (e.g., \texttt{(G := G)}).
	\item \textbf{Equivalence mechanics:} Constructing \texttt{sigmaPerm} as an \texttt{Equiv.Perm X}
	      required providing both \texttt{left\_inv} and \texttt{right\_inv} proofs, which meant carefully
	      applying the \texttt{ga\_mul} axiom in both directions.
	\item \textbf{Coercions:} Working with subgroups required understanding how Lean coerces elements
	      of type \texttt{H} (a subgroup) to type \texttt{G} (the ambient group) using \texttt{(h : G)}.
	\item \textbf{Function extensionality:} Proving equality of permutations required \texttt{Equiv.ext},
	      and equality of vector-valued functions required \texttt{ext i}, which were not immediately obvious.
\end{itemize}
\iffalse
	\subsection{Potential extensions}
	This development could be extended in several directions:
	\begin{itemize}
		\item \textbf{Orbit-stabilizer theorem:} Formalize the relationship $|G| = |\text{Orb}(x)| \cdot |G_x|$
		      for finite groups, which requires defining orbits and proving Lagrange's theorem.
		\item \textbf{Burnside's lemma:} Count orbits using the formula involving fixed points.
		\item \textbf{Actions on cosets:} Show that a group acts on the left cosets of a subgroup,
		      yielding the standard permutation representation.
		\item \textbf{Cayley's theorem:} Specialize the permutation representation to the case where
		      $G$ acts on itself, proving every group embeds in a symmetric group.
	\end{itemize}
\fi
\section{AI Assistance Statement}
I used AI-assisted tools in a limited and strictly auxiliary manner. Specifically, AI was used to help automate the conversion of Lean source files into \LaTeX\ code blocks for inclusion in this report, and to assist with reorganising a single Lean file into multiple smaller files during the development process. In addition, AI was occasionally used for high-level brainstorming about structure and presentation.

All mathematical content, Lean definitions, proofs, and formal statements were written, checked, and verified by me. AI tools were not used to replace mathematical reasoning or to generate unverified results. I take full responsibility for the correctness and originality of both the Lean code and the written exposition.


\section{References}
John B. Fraleigh, Victor J. Katz, \textit{A First Course in Abstract Algebra},
Addison--Wesley, 2003, Section 16 (Group Actions).

\end{document}
